%-------------------------------------------------------------------------------
% Summary
%-------------------------------------------------------------------------------

\chapter*{Summary}

\addcontentsline{toc}{chapter}{Summary}

Continuous black-box optimization plays an important role in parameter extraction and complex engineering simulations, yet traditional human-engineered metaheuristics can perform poorly when operating under noisy evaluations. In real-world applications, empirical evaluations are frequently affected by stochastic noise that obscures the underlying optimization landscape and can mislead conventional gradient-free search methods. Addressing these challenges has traditionally required human experts to develop specialized mechanisms, such as filtering or step-size adaptation rules, through extensive manual experimentation, making the process labor-intensive and difficult to scale across structurally diverse optimization problems. To address this limitation, this thesis investigates Automated Algorithm Design (AAD) using Large Language Models (LLMs), moving beyond static parameter tuning toward the autonomous synthesis of executable optimization algorithms. Building on the Large Language Model Evolutionary Algorithm (LLaMEA) framework, the proposed approach investigates whether evolutionary feedback obtained from noisy fitness evaluations can lead LLMs to synthesize optimization heuristics that are better suited to stochastic conditions. The study uses controlled noise applied to continuous black-box benchmark functions to provide a repeatable evaluation setting. Algorithms evolved under noisy and noiseless conditions will be evaluated across different noise settings, including clean evaluations, and compared in terms of performance, robustness, and generalization. Established optimization methods may also be included as reference baselines. This evaluation will provide an empirical assessment of how exposure to noisy fitness evaluations during algorithm evolution affects the resulting optimization heuristics.
